% ID: SEL-D0003-OB-E04
% ESTADO: validada. S004, 2026-11-07; tres pasadas del mismo autor y revisión matemática, didáctica y visual.
\textbf{a) Intervalo para la proporción poblacional.}\par\nopagebreak[4]
La estimación es $\widehat p=2250/5000=0.45$. Suponemos muestreo aleatorio representativo con observaciones independientes; no hay tamaño poblacional para una corrección por población finita. Los recuentos 2250 y 2750 permiten la aproximación normal usual.

Para confianza 97.5\,\%, $\alpha=0.025$ y cada cola tiene 0.0125, de modo que $\Phi(z)=0.9875$ y $z\approx2.241403$. El intervalo aproximado es
\[
\widehat p\pm z\sqrt{\frac{\widehat p(1-\widehat p)}n}.
\]
\[
E=2.241403\sqrt{\frac{0.45(0.55)}{5000}}\approx0.015769666.
\]
\[
\boxed{IC_{\text{97.5\%}}\approx(0.434230334,\ 0.465769666).}
\]
En porcentajes, aproximadamente entre 43.42\,\% y 46.58\,\%. La tabla da $\Phi(2.24)=0.9875$ redondeado; usar 2.24 ofrece un intervalo equivalente a la precisión porcentual indicada. La confianza corresponde a la cobertura del método en muestreos repetidos.

\textbf{b) Nueva proporción y nueva confianza.}\par\nopagebreak[4]
Ahora se usan $\widehat p=0.5$ y confianza 92.5\,\%, no los valores del apartado anterior. Tenemos $\alpha=0.075$, $1-\alpha/2=0.9625$ y $z_{0.9625}\approx1.780464342$. El error debe ser estrictamente menor que 0.03:
\[
z_{0.9625}\sqrt{\frac{0.5(0.5)}n}<0.03
\iff n>\frac{z_{0.9625}^2(0.25)}{0.03^2}
\approx880.570353.
\]
El tamaño total mínimo es \boxed{n=881\text{ estudiantes}}. Verificación:
\[
E(880)\approx0.030009720>0.03,\qquad
E(881)\approx0.029992684<0.03.
\]
La tabla da $\Phi(1.78)=0.9625$ redondeado; con 1.78 también resulta 881. El error 0.03 son tres puntos porcentuales, no el 3\,\% de la proporción 0.5.
